\begin{table}[!ht]
\centering
\caption{Dataset and ground-truth summary.}
\label{tab:dataset_gt_summary}
\small
\begin{tabular}{lccccc}
\toprule
\textbf{Scope} & \textbf{Docs} & \textbf{Req. items} & \textbf{Dimensions} & \textbf{Units} & \textbf{GT positives} \\
\midrule
PURE subset & 10 & 978 & U/A/C/V & 3,912 & 240 (6.13\%) \\
\bottomrule
\end{tabular}

\vspace{0.6em}

\begin{tabular}{lccccc}
\toprule
\textbf{GT distribution} & \textbf{U} & \textbf{A} & \textbf{C} & \textbf{V} & \textbf{Total} \\
\midrule
Positive units & 44 & 135 & 4 & 57 & 240 \\
Positive rate & 4.50\% & 13.80\% & 0.41\% & 5.83\% & 6.13\% \\
\bottomrule
\end{tabular}

\par\vspace{0.4em}
{\footnotesize \textit{Note.} The evaluation unit is an $(item,type)$ pair, where $type \in \{U,A,C,V\}$ denotes Understandable, Unambiguous, Correctness, and Verifiable, respectively. Ground-truth labels are binary: $1$ indicates a quality issue in the corresponding dimension, $0$ otherwise. The final ground truth was built from two independent expert annotations, with disagreements resolved by a third senior adjudicator.\par}
\end{table}
